\section{Justificación}

La transcripción automática de la música presenta un desafío al momento de procesar las señales digitales, de igual forma constituye un problema que toma gran relevancia al momento de recuperar información musical (MIR), esta tiene un gran potencial para transformar la manera en que se analiza, se interpreta y se accede al contenido musical~\cite{jamshidi2024machine}. A pesar de los avances recientes en los ámbitos de machine learning, lo sistemas actuales no tienen la capacidad máxima de precisión y flexibilidad al comparar con un músico profesional, ya que estos tienen la capacidad de entender información polifónica mucho más amplia, en estos escenarios polifónicos complejos se convergen múltiples voces, sonidos y narrativas sonoras que son complejas para una maquina de entender de manera simultanea~\cite{benetos2019automatic}.

El desarrollo de aplicaciones que permitan la extracción de información como acordes, estructuras musicales, datos de frecuencia y demás, puede llegar a ser una herramienta clave para un aprendizaje musical, que puede acompañar al interesado y así darle un respectivo uso educativo. La base para el desarrollo de una aplicación parte de una transcripción musical, donde permite la conversión de información auditiva en representaciones simbólicas, esto constituye la base de un estudio, una interpretación y una aplicación en respectivos contextos~\cite{jamshidi2024machine}. En este sentido, automatizar este proceso contribuye en una reducción del tiempo y de dificultad a interpretadores cuya tarea sea transcribir de forma manual y presentarle más información que no sea solo en formato de partitura.

